\section{Interlocking Service Module}
\label{sec:methodology:interlocking-service}

The InterlockingService module acts as the central coordinator between operator or system-initiated requests and the underlying safety branches. It governs the lifecycle of interlocking operations, beginning with the initialization of validation logic, through rule execution, and into reactive safety enforcement. By supervising dedicated branches for signals, point machines, and track circuits, the service consolidates their outputs into a single validation result, ensuring consistency across the system. This result not only determines whether a route or signal request may proceed but also provides reason-tagged feedback for both the operator interface and the system logger, reinforcing transparency and accountability.

Inputs to the service may originate from a variety of sources, including the operator GUI, automated testing scripts, or the simulation environment. Once received, these requests are interpreted against current infrastructure states retrieved from the database, after which the service coordinates validation checks across the relevant modules. The outcome is then disseminated in real time: updates are reflected on the graphical simulation panel, logs are recorded for traceability, and, where necessary, alerts or freeze events are triggered to contain unsafe conditions.

The design emphasizes modularity, responsiveness, and resilience. Validation logic is branch-based, ensuring that point, track, and signal rules remain cleanly separated and individually testable. Performance is benchmarked so that even under heavy operational loads, the service can respond within real-time constraints. Reactive safety is also embedded: changes in track occupancy immediately trigger enforcement logic, enabling the system to respond dynamically to unexpected events. In this role, the InterlockingService functions as the orchestration layer of the interlocking system, binding together the rule engine, validation modules, and operator interface into a coherent and fail-safe whole.